% ----------------------------------------------------------
% Introdução
% ----------------------------------------------------------
\chapter{Introdução} \label{cha:introducao}

A indústria de jogos digitais se configura como um dos segmentos do entretenimento que mais crescem e faturam no mundo, estima-se que ele tenha movimentado cerca de US\$ 188,8 bilhões em 2025, com um público global superior a três bilhões de jogadores \cite{newzoo2025}. Paralelamente, a maneira como os jogos ocupam espaço no mercado tem se transformado, pois o sucesso de um título passou a depender principalmente da formação de comunidades e do conteúdo criado pelos própios jogadores \cite{bain2025}.

Nesse contexto, ouvir a comunidade se tornou parte importante no ciclo de desenvolvimento, principalmente para estúdios independentes, que veem nas contribuições da comunidade um meio acessível para embasar decisões criativas e técnicas durante o projeto \cite{uchoa2018}. O Desenvolvimento Aberto (\textit{Open Development}), apresentado por \citeonline{thominet2018}, aprofunda esssa ideia ao propor que o público seja envolvido de forma contínua, testando versões e tendo seu retorno incorporado ao produto.

Contudo, atualmente o retorno da comunidade raramente fica concentrado em um único lugar, Ele se distribui por diferentes plataformas, como fóruns, redes sociais e plataformas de distribuição digital, além dessa fragmentação, cada ambiente possui seu próprio tipo de conteúdo, que pode variar de acordo com o público e a forma como ele interage com o jogo \cite{tong2022}. Essa dispersão de informações dificulta a análise do retorno da comunidade, tornando mais custoso lidar com o crescente volume de dados \cite{xiao2025}. 

Apesar da existência de canais em que os desenvolvedores podem interagir com a comunidade, lidar com essa multiplicidade de plataformas e conteúdos ainda é um problema, pois pode resultar em uma experiência fragmentada e confusa para os usuários. Além disso, a ausência de uma interface unificada pode dificultar a análise e a priorização dos \textit{feedbacks}, tornando mais desafiador identificar padrões, tendências e pontos de melhoria.

Sendo assim, a usabilidade da interface torna-se um aspecto importante para o desenvolvimento de uma plataforma que centralize essas interações. Segundo \citeonline{nielsen1993}, a usabilidade da interface de um sistema desempenha um papel essencial na aceitabilidade prática do sistema interativo. Em ambientes de uso voluntário e comunitário, a ausência de barreiras de interação e uma experiência positiva são determinantes para evitar a frustração do público, garantindo a adesão e o engajamento contínuo dos usuários.

Diante desse cenário, este trabalho propõe a prototipação e o desenvolvimento de uma interface que centralize esses conteúdos, com o objetivo de facilitar a comunicação entre jogadores e desenvolvedores e oferecer uma experiência mais fluida e eficiente aos usuários. A interface será validada por meio de testes com usuários, afim de verificar a adequação do layout, da navegação e das principais funcionalidades.


% Diante desse cenário, este trabalho propõe a prototipação e o desenvolvimento de uma interface que centralize esses conteúdos, com o objetivo de facilitar a comunicação entre jogadores e desenvolvedores e oferecer uma experiência mais fluida e eficiente aos usuários.
